\everymath{\displaystyle}

%%%%%%%%%%%%%%%%%%%%%%%%%%%%%%%%%%%%%%%%
%           main body
%%%%%%%%%%%%%%%%%%%%%%%%%%%%%%%%%%%%%%%%

%-----------------------------------
%	TITLE PAGE
%-----------------------------------

\title[第十章\\
函数项级数]{\texorpdfstring{\LARGE \bfseries 第十章\quad 函数项级数}{第十章 函数项级数}} % The short title appears at the bottom of every slide, the full title is only on the title page
\author[\smaller \S 10.1\\
一致收敛性] % Your institution as it will appear on the bottom of every slide, maybe shorthand to save space
{\Large \bfseries \S 10.1 函数项级数的一致收敛性}
% \author{Singular Value Decomposition} % Your name
% \date{\today} % Date, can be changed to a custom date
\date{}

%------------------------------------------------

\begin{document}

%------------------------------------------------

\begin{frame}

\titlepage % Print the title page as the first slide

\end{frame}

%-----------------------------------
%	PRESENTATION SLIDES
%-----------------------------------

%------------------------------------------------

\begin{frame}
\frametitle{绪言}

设有一列函数 $\{u_n(x)\},$ 其公共定义域为 $E.$ 那么对每个固定的 $x_0 \in E,$ $\{u_n(x_0)\}$ 就是一个数列. 可以考虑相应的数项级数
\[
\sum_{n=1}^\infty u_n(x_0).
\]
将所有使得上述数项级数收敛的 $x_0$ 的集合记为 $D,$ 称为\alert{收敛域}, 即
\[
D = \{x \in E: \sum_{n=1}^\infty u_n(x) \text{ 收敛}\}.
\]
\pause
可以在 $D$ 上定义一个新的函数
\[
\sum_{n=1}^\infty u_n: D \to \mathbb{R}, \quad x \mapsto \sum_{n=1}^\infty u_n(x).
\]
$\sum_{n=1}^\infty u_n$ 就是所谓的\alert{函数项级数}, 即它的通项为函数.

\end{frame}

%------------------------------------------------

\begin{frame}
\frametitle{绪言}

{
\larger[1]

有时, \Circled{1} $\sum_{n=1}^\infty u_n$ 的收敛性``不够好'', 例如与积分, 微分等运算的不可交换,
\pause
这时, 我们需要将其条件加强, 要求收敛性满足某些``均匀''的条件, 例如一致收敛;

\pause
\vspace{2.5em}

有时, \Circled{2} 需要将可研究的范围扩大, 考虑更弱一些意义下的收敛性,
\pause
例如, 依测度收敛, 依范数收敛, 弱收敛, 等等. 这是后续课程的内容了.
}

\end{frame}

%------------------------------------------------

\section{函数项级数的点态收敛}

%------------------------------------------------

\begin{frame}
\frametitle{函数项级数的基本概念}

\begin{Def}[函数项级数]
设 $\{u_n(x)\}$ 是定义在公共定义域 $E$ 上的一列函数. 形式地将他们的``和''
\[
u_1(x) + u_2(x) + \cdots + u_n(x) + \cdots
\]
称为\alert{函数项级数}, 记为 $\sum_{n=1}^\infty u_n(x).$
\end{Def}

\pause

\begin{Def}[函数项级数的收敛域与和函数]
设 $x_0 \in E.$ 如果数项级数 $\sum_{n=1}^\infty u_n(x_0)$ 收敛, 则称函数项级数 $\sum_{n=1}^\infty u_n(x)$ 在 $x_0$ 处收敛, 或者说 $x_0$ 是函数项级数 $\sum_{n=1}^\infty u_n(x)$ 的一个收敛点. 函数项级数 $\sum_{n=1}^\infty u_n(x)$ 的所有收敛点的集合 $D$ 称为它的\alert{收敛域}. 在 $D$ 上定义函数 $S(x) = \sum_{n=1}^\infty u_n(x), x \in D,$ 则称 $S$ 是函数项级数 $\sum_{n=1}^\infty u_n$ 的\alert{和函数}.
\end{Def}

\end{frame}

%------------------------------------------------

\begin{frame}
\frametitle{函数项级数的基本概念}

\begin{remark}
由于以上定义的和函数 $S$ 是通过逐点定义的方式得到的, 因此, 称函数项级数 $\sum_{n=1}^\infty u_n$ 在 $D$ 上\alert{点态收敛}于 $S.$
\end{remark}

\pause

\begin{eg}
\begin{enumerate}
\item $\sum_{n=1}^\infty x^n$ 的收敛域为 $(-1, 1)$, 和函数为 $S(x) = \frac{1}{1-x}.$
\pause
\item $\sum_{n=1}^\infty \frac{x^n}{n}$ 的收敛域为 $(-1, 1) \cup \{-1\} = [-1, 1)$, 和函数为 $S(x) = -\ln(1-x).$
\pause
注意, 该函数定义域是严格包含收敛域的!
\pause
\item $\sum_{n=1}^\infty \frac{x^n}{n!}$ 的收敛域为 $(-\infty, +\infty)$, 和函数为 $S(x) = e^x - 1.$
\end{enumerate}
\end{eg}

\end{frame}

%------------------------------------------------

\begin{frame}
\frametitle{函数项级数的部分和函数}

类似于数项级数, 可以考察函数项级数 $\sum_{n=1}^\infty u_n$ 的\alert{部分和}, 即
\begin{Def}[函数项级数的部分和函数]
对函数项级数 $\sum_{n=1}^\infty u_n,$ 定义其\alert{部分和函数, 或者说前 $n$ 项和函数} 为
\[
S_n(x) = \sum_{k=1}^n u_k(x), \quad n = 1, 2, \cdots.
\]
\end{Def}

\pause

\begin{remark}

\end{remark}

\end{frame}

%------------------------------------------------

\end{document}
